\documentclass[border=0pt]{standalone}
% The shared preamble of every docs diagram (docs/diagrams/<name>.tex).
% dvisvgm's driver first, so patterns and opacity survive into the SVG.
\def\pgfsysdriver{pgfsys-dvisvgm.def}
\usepackage[siunitx,RPvoltages]{circuitikz}
\usepackage{xcolor}
% Latin Modern: outline fonts for every glyph a label uses, µ included.
\usepackage[T1]{fontenc}
\usepackage{lmodern}
\usetikzlibrary{arrows.meta,backgrounds,calc,patterns,positioning}
\tikzset{
  % A white card, so diagrams read in mdBook's dark themes too.
  every picture/.append style={show background rectangle,
    background rectangle/.style={fill=white,inner sep=4pt}},
  % Layout geometry: layers told apart by fill and hatching, never colour alone.
  layer a/.style={draw=black,line width=0.6pt,fill=gray!20},
  layer b/.style={draw=black,line width=0.6pt,fill=gray!45},
  layer c/.style={draw=black,line width=0.6pt,fill=white,postaction={pattern=north east lines}},
  layer d/.style={draw=black,line width=0.6pt,fill=white,postaction={pattern=crosshatch dots}},
  marker/.style={draw=black,dashed,line width=0.6pt},
  % A measured distance and its label.
  dim/.style={{Latex[length=2mm]}-{Latex[length=2mm]},line width=0.5pt},
  dim label/.style={font=\small,fill=white,inner sep=1pt},
  % Where a rule fails: thick, dashed and labelled, not only red.
  violation/.style={draw=red!75!black,line width=1.4pt,dash pattern=on 3pt off 2pt},
  violation label/.style={font=\small\bfseries,text=red!75!black},
  caption/.style={font=\small},
}

\begin{document}
\begin{circuitikz}
  % Two inverters drive Y1 and Y2; one nfet joins them.
  \node[not port] (i1) at (0,0) {};
  \node[not port,rotate=180] (i2) at (7,0) {};
  \draw (i1.in) -- ++(-0.4,0) node[left]{A};
  \draw (i2.in) -- ++(0.4,0) node[right]{B};
  \node[caption,anchor=south] at (0,0.75) {VDD, VSS};
  \node[caption,anchor=south] at (7,0.75) {VDD, VSS};
  \draw (i1.out) to[short,-*] (2.2,0) coordinate(y1);
  \draw (i2.out) to[short,-*] (4.8,0) coordinate(y2);
  \node[caption,anchor=south] at (y1) {Y1};
  \node[caption,anchor=south] at (y2) {Y2};
  \draw (2.7,0) node[nmos,rotate=-90,anchor=S](t){};
  \draw (y1) -- (t.S); \draw (t.D) -- (y2);
  \draw (t.G) -- ++(0,0.4) node[above]{EN};
  \draw[violation] ($(t.S)+(0,-0.75)$) rectangle ($(t.D)+(0,0.95)$);
  \node[violation label,anchor=north] at ($(t.S)!0.5!(t.D)+(0,-0.85)$) {L = 200 nm $\leq$ 200 nm, used both ways};
  \node[caption,anchor=north,align=center] at ($(t.S)!0.5!(t.D)+(0,-1.35)$) {cut out, Y1 and Y2 each still reach VDD and VSS};
\end{circuitikz}
\end{document}
